\section{Summary}
The purpose of this work is to investigate the impacts of 
deploying \gls{HALEU}-fueled advanced reactors in the US. The lit 
review discussed the tools used to perform this investigation, 
and the current understanding of the fuel cycle and neutronics 
effects of using \gls{HALEU} in reactors. 

The results presented thus far in the investigation show 

\section{Propsed Work}
\subsection{Recycle Fuel Cycles}
The once-through fuel cycles investigated will be modified to include recycling 
to investigate how this type of \gls{NFC} affects the material requirements 
at each time, and how the recycling scheme impacts them as well. 
The scenarios investigated for this section of analysis vary by the 
energy demand of the scenario and the recycling scheme. Energy demands vary 
between a no growth and a 1\% growth model, the same as how this parameter 
is varied in the once-through fuel cycle models. The reycling scheme is varied 
to be either a limited or a continuous recycle, similar to how this parameter 
was varied in the \gls{ES} \cite{wigeland_nuclear_2014}. All of the scenarios 
with recycling will model the transition from the current fleet of \glspl{LWR} 
to the X-energy Xe-100, \gls{USNC} \gls{MMR}, and the NuScale VOYGR, the same 
as in scenarios 7 and 13. 

\begin{table}[ht]
    \centering
    \caption{Summary of the recycle fuel cycle transition scenarios.}
    \label{tab:scenarios_recycle}
    \begin{tabular}{l l l}
            \hline
            Scenario number & Energy growth model & Recycle scheme\\\hline
            14 & No growth & Limited \\
            15 & No growth & Continuous \\
            16 & 1\% growth & Limited \\
            17 & 1\% growth & Continuous \\
            \hline
    \end{tabular}
\end{table}

Flow of material through the fuel cycle scenarios with recycling is shown in 
Figure \ref{fig:recycle_fuel_cycle}. This fuel cycle assumes that \gls{SNF}
from the \glspl{LWR} and advanced reactors can undergo separation and be 
recycled.

\begin{figure}
    \centering
    \begin{tikzpicture}[node distance=1.5cm]
        \node (mine) [facility] {Uranium Mine};
        \node (enrichment) [facility, below of=mine]{Enrichment};
        \node (reactor) [facility, below of=enrichment]{Reactor};
        \node (adv_reactor) [transition, right of=reactor, xshift=3cm]{Advanced Reactor};
        \node (wetstorage) [facility, below of=reactor]{Wet Storage};
        \node (drystorage) [facility, below of=wetstorage]{Dry Storage};
        \node (cooling) [transition, below of=adv_reactor]{Cooling Pool};
        \node (sinkhlw) [facility, below of=drystorage, xshift=2.5cm]{HLW Sink};
        \node (sinkllw) [facility, left of=enrichment, xshift=-3cm]{LLW Sink};
        \node (separation) [transition, below of=cooling]{Separations};
        \node (fuelfab) [transition, below of=adv_reactor,xshift=3cm]{Fuel Fab};

        \draw [arrow] (mine) -- node[anchor=east]{Natural U} (enrichment);
        \draw [arrow] (enrichment) -- node[anchor=east]{Enriched U}(reactor);
        \draw [arrow] (enrichment) -- node[anchor=south]{Tails}(sinkllw);
        \draw [arrow] (enrichment) -| node[anchor=west]{Fresh Fuel}(adv_reactor);
        \draw [arrow] (reactor) -- node[anchor=east]{Spent UOX}(wetstorage);
        \draw [arrow] (wetstorage) -- node[anchor=east]{Cool Spent UOX}(drystorage);
        \draw [arrow] (drystorage) |- node[anchor=east]{Casked Spent UOX}(sinkhlw);
        \draw [arrow] (adv_reactor) -- node[anchor=east]{Spent Fuel}(cooling);
        \draw [arrow] (cooling) -- node[anchor=east]{Cooled Spent Fuel}(separation);
        \draw [arrow] (separation) -| node[anchor=west]{Separated fissile material}(fuelfab);
        \draw [arrow] (fuelfab) |- node[anchor=south]{Reprocessed fuel}(adv_reactor);
        \draw [arrow] (separation) |- node[anchor=west]{Fission Products}(sinkhlw);

        \end{tikzpicture}
    \caption{Fuel cycle facilities and material flow between facilities. Facilities in 
    red are added in for the transition scenarios.}
    \label{fig:recycle_fuel_cycle}
\end{figure}

A separations facility is defined to be deployed 5 years before the transition 
begins (i.e. in 2020), as this strategy is commonly used for modeling 
transitions with reyccling to ensure that enough fuel can be spearated and 
processed in time for use in advanced reactors 
\cite{passerini_systematic_2014,richards_application_2021}.
Assume processing losses of 1\% \cite{wigeland_nuclear_2014}.

\subsection{Sensitivity Analysis and Optimization}
The sensitivity analysis is designed to identify how 
the variation of select hyperparameters affect select performance 
metrics of the fuel cycle and how the hyperparameters synergistically 
affect the performance metrics. This information is then used to design 
a transition scenario that is optimized for each of the performance metrics. 

To perform the sensitivity analysis, uncertainty quantification, and 
transition optimization \Cyclus was coupled to Dakota \cite{adams_dakota_2019},
and Dakota was used as the driver for the analysis. The metrics of 
interest for this analysis include the amount of waste generated that 
must be sent to a repository, the mass of natural uranium required, and 
the amount of \gls{SWU} required. The first two performance 
metrics used in this work match metrics considered in the \gls{ES} 
\cite{wigeland_nuclear_2014}, and all three of these metrics were considered 
in previous sensitivity analysis when Dakota was coupled to \gls{DYMOND}
\cite{feng_sensitivity_2020,richards_application_2021}. 

\subsection{Neutronics Analysis}
Recovery and downblending of \gls{HALEU} to produce \gls{HALEU} means that 
the fuel is likely to contain an impurities that are present in the 
\gls{HEU}. Known impurities in potential \gls{HEU}
supplies to create \gls{HALEU} include uranium-232 and uranium-236
\cite{vaden_isotopic_2018,nelson_foreign_2010},  
which are parasitic neutron absorbers. This work models the use of 
impure \gls{HALEU} fuel in a reactor and compares it to \gls{HALEU} fuel 
without these impurities to understand how the impurities affect 
the performance of the reactor. 

Full core models of the Xe-100 and \gls{MMR} will be created in Serpent 
\cite{leppanen_serpent_2014} to model the neutronics of each reactor 
design. The fuel composition for each core will be varied to provide 
a comparison of pure (containing only U-235 and U-238) and impure fuel.
Material compositions of impure fuel will be modeled based on both the 
known compositions of \gls{HALEU} from the \gls{EBR} stockpile 
\cite{vaden_isotopic_2018} and the Y-12 stockpile \cite{nelson_foreign_2010}.
Each core will be modeled three times, for each fuel composition considered.
By modeling the core composition to be entirely composed of impure fuel, 
this work reports the most extreme effect that these impurities can 
have on reactor performance.  

Results of each simulation that will be compared include the neutron flux at 
\gls{BOL}, mid-cycle, and \gls{EOL}, $k_{eff}$ as a function of burnup, 
fuel temperature reactivity coefficient, and 
$\beta_{eff}$. Each of these parameters has impacts on performance of 
the reactor, such as the materials degredation rate, amount of burnable 
poisons required, control rod worth, and the cycle time. Investigating 
each of these results will help to determine if the impurities potentially 
present in \gls{HALEU} will prevent any of the design criteria of the 
reactors from being met. 